% Options for packages loaded elsewhere
% Options for packages loaded elsewhere
\PassOptionsToPackage{unicode}{hyperref}
\PassOptionsToPackage{hyphens}{url}
%
\documentclass[
  11pt,
  letterpaper,
  DIV=11,
  numbers=noendperiod,
  twoside]{scrartcl}
\usepackage{xcolor}
\usepackage[left=1.1in, right=1in, top=0.8in, bottom=0.8in,
paperheight=9.5in, paperwidth=7in, includemp=TRUE, marginparwidth=0in,
marginparsep=0in]{geometry}
\usepackage{amsmath,amssymb}
\setcounter{secnumdepth}{3}
\usepackage{iftex}
\ifPDFTeX
  \usepackage[T1]{fontenc}
  \usepackage[utf8]{inputenc}
  \usepackage{textcomp} % provide euro and other symbols
\else % if luatex or xetex
  \usepackage{unicode-math} % this also loads fontspec
  \defaultfontfeatures{Scale=MatchLowercase}
  \defaultfontfeatures[\rmfamily]{Ligatures=TeX,Scale=1}
\fi
\usepackage{lmodern}
\ifPDFTeX\else
  % xetex/luatex font selection
  \setmathfont[]{Garamond-Math}
\fi
% Use upquote if available, for straight quotes in verbatim environments
\IfFileExists{upquote.sty}{\usepackage{upquote}}{}
\IfFileExists{microtype.sty}{% use microtype if available
  \usepackage[]{microtype}
  \UseMicrotypeSet[protrusion]{basicmath} % disable protrusion for tt fonts
}{}
\usepackage{setspace}
% Make \paragraph and \subparagraph free-standing
\makeatletter
\ifx\paragraph\undefined\else
  \let\oldparagraph\paragraph
  \renewcommand{\paragraph}{
    \@ifstar
      \xxxParagraphStar
      \xxxParagraphNoStar
  }
  \newcommand{\xxxParagraphStar}[1]{\oldparagraph*{#1}\mbox{}}
  \newcommand{\xxxParagraphNoStar}[1]{\oldparagraph{#1}\mbox{}}
\fi
\ifx\subparagraph\undefined\else
  \let\oldsubparagraph\subparagraph
  \renewcommand{\subparagraph}{
    \@ifstar
      \xxxSubParagraphStar
      \xxxSubParagraphNoStar
  }
  \newcommand{\xxxSubParagraphStar}[1]{\oldsubparagraph*{#1}\mbox{}}
  \newcommand{\xxxSubParagraphNoStar}[1]{\oldsubparagraph{#1}\mbox{}}
\fi
\makeatother


\usepackage{longtable,booktabs,array}
\newcounter{none} % for unnumbered tables
\usepackage{calc} % for calculating minipage widths
% Correct order of tables after \paragraph or \subparagraph
\usepackage{etoolbox}
\makeatletter
\patchcmd\longtable{\par}{\if@noskipsec\mbox{}\fi\par}{}{}
\makeatother
% Allow footnotes in longtable head/foot
\IfFileExists{footnotehyper.sty}{\usepackage{footnotehyper}}{\usepackage{footnote}}
\makesavenoteenv{longtable}
\usepackage{graphicx}
\makeatletter
\newsavebox\pandoc@box
\newcommand*\pandocbounded[1]{% scales image to fit in text height/width
  \sbox\pandoc@box{#1}%
  \Gscale@div\@tempa{\textheight}{\dimexpr\ht\pandoc@box+\dp\pandoc@box\relax}%
  \Gscale@div\@tempb{\linewidth}{\wd\pandoc@box}%
  \ifdim\@tempb\p@<\@tempa\p@\let\@tempa\@tempb\fi% select the smaller of both
  \ifdim\@tempa\p@<\p@\scalebox{\@tempa}{\usebox\pandoc@box}%
  \else\usebox{\pandoc@box}%
  \fi%
}
% Set default figure placement to htbp
\def\fps@figure{htbp}
\makeatother


% definitions for citeproc citations
\NewDocumentCommand\citeproctext{}{}
\NewDocumentCommand\citeproc{mm}{%
  \begingroup\def\citeproctext{#2}\cite{#1}\endgroup}
\makeatletter
 % allow citations to break across lines
 \let\@cite@ofmt\@firstofone
 % avoid brackets around text for \cite:
 \def\@biblabel#1{}
 \def\@cite#1#2{{#1\if@tempswa , #2\fi}}
\makeatother
\newlength{\cslhangindent}
\setlength{\cslhangindent}{1.5em}
\newlength{\csllabelwidth}
\setlength{\csllabelwidth}{3em}
\newenvironment{CSLReferences}[2] % #1 hanging-indent, #2 entry-spacing
 {\begin{list}{}{%
  \setlength{\itemindent}{0pt}
  \setlength{\leftmargin}{0pt}
  \setlength{\parsep}{0pt}
  % turn on hanging indent if param 1 is 1
  \ifodd #1
   \setlength{\leftmargin}{\cslhangindent}
   \setlength{\itemindent}{-1\cslhangindent}
  \fi
  % set entry spacing
  \setlength{\itemsep}{#2\baselineskip}}}
 {\end{list}}
\usepackage{calc}
\newcommand{\CSLBlock}[1]{\hfill\break\parbox[t]{\linewidth}{\strut\ignorespaces#1\strut}}
\newcommand{\CSLLeftMargin}[1]{\parbox[t]{\csllabelwidth}{\strut#1\strut}}
\newcommand{\CSLRightInline}[1]{\parbox[t]{\linewidth - \csllabelwidth}{\strut#1\strut}}
\newcommand{\CSLIndent}[1]{\hspace{\cslhangindent}#1}



\setlength{\emergencystretch}{3em} % prevent overfull lines

\providecommand{\tightlist}{%
  \setlength{\itemsep}{0pt}\setlength{\parskip}{0pt}}



 


\setlength\heavyrulewidth{0ex}
\setlength\lightrulewidth{0ex}
\usepackage[automark]{scrlayer-scrpage}
\clearpairofpagestyles
\cehead{
  Claude
  }
\cohead{
  A Counterexample to Arrhenius’s Sixth Impossibility Theorem
  }
\ohead{\bfseries \pagemark}
\cfoot{}
\makeatletter
\newcommand*\NoIndentAfterEnv[1]{%
  \AfterEndEnvironment{#1}{\par\@afterindentfalse\@afterheading}}
\makeatother
\NoIndentAfterEnv{itemize}
\NoIndentAfterEnv{enumerate}
\NoIndentAfterEnv{description}
\NoIndentAfterEnv{quote}
\NoIndentAfterEnv{equation}
\NoIndentAfterEnv{longtable}
\NoIndentAfterEnv{abstract}
\renewenvironment{abstract}
 {\vspace{-1.25cm}
 \quotation\small\noindent\emph{Abstract}:}
 {\endquotation}
\newfontfamily\tfont{EB Garamond}
\addtokomafont{disposition}{\rmfamily}
\addtokomafont{descriptionlabel}{\rmfamily}
\addtokomafont{title}{\normalfont\itshape}
\let\footnoterule\relax

\makeatletter
\renewcommand{\@maketitle}{%
  \newpage
  \null
  \vskip 2em%
  \begin{center}%
  \let \footnote \thanks
    {\itshape\huge\@title \par}%
    \vskip 0.5em%  % Reduced from default
    {\large
      \lineskip 0.3em%  % Reduced from default 0.5em
      \begin{tabular}[t]{c}%
        \@author
      \end{tabular}\par}%
    \vskip 0.5em%  % Reduced from default
    {\large \@date}%
  \end{center}%
  \par
  }
\makeatother
\RequirePackage{lettrine}

\renewenvironment{abstract}
 {\quotation\small\noindent\emph{Abstract}:}
 {\endquotation\vspace{-0.02cm}}

\setmainfont{EB Garamond Math}[
  BoldFont = {EB Garamond SemiBold},
  ItalicFont = {EB Garamond Italic},
  RawFeature = {+smcp},
]

\newfontfamily\scfont{EB Garamond Regular}[RawFeature=+smcp]
\renewcommand{\textsc}[1]{{\scfont #1}}

\renewcommand{\LettrineTextFont}{\scfont}
% Number theorem-like environments straight through (Theorem 1, Theorem 2, ...)
% rather than by section (1.1, 1.2, ...), which is Quarto's LaTeX default.
% Guarded, since Quarto only defines a counter for environments the document uses.
\makeatletter
\AtBeginDocument{%
  \@for\bw@thm:={theorem,lemma,corollary,proposition,conjecture,definition,%
                 example,exercise,refremark,refsolution}\do{%
    \@ifundefined{c@\bw@thm}{}{\expandafter\counterwithout\expandafter{\bw@thm}{section}}%
  }%
}
\makeatother
\KOMAoption{captions}{tableheading}
\makeatletter
\@ifpackageloaded{caption}{}{\usepackage{caption}}
\AtBeginDocument{%
\ifdefined\contentsname
  \renewcommand*\contentsname{Table of contents}
\else
  \newcommand\contentsname{Table of contents}
\fi
\ifdefined\listfigurename
  \renewcommand*\listfigurename{List of Figures}
\else
  \newcommand\listfigurename{List of Figures}
\fi
\ifdefined\listtablename
  \renewcommand*\listtablename{List of Tables}
\else
  \newcommand\listtablename{List of Tables}
\fi
\ifdefined\figurename
  \renewcommand*\figurename{Figure}
\else
  \newcommand\figurename{Figure}
\fi
\ifdefined\tablename
  \renewcommand*\tablename{Table}
\else
  \newcommand\tablename{Table}
\fi
}
\@ifpackageloaded{float}{}{\usepackage{float}}
\floatstyle{ruled}
\@ifundefined{c@chapter}{\newfloat{codelisting}{h}{lop}}{\newfloat{codelisting}{h}{lop}[chapter]}
\floatname{codelisting}{Listing}
\newcommand*\listoflistings{\listof{codelisting}{List of Listings}}
\makeatother
\makeatletter
\makeatother
\makeatletter
\@ifpackageloaded{caption}{}{\usepackage{caption}}
\@ifpackageloaded{subcaption}{}{\usepackage{subcaption}}
\makeatother
\usepackage{bookmark}
\IfFileExists{xurl.sty}{\usepackage{xurl}}{} % add URL line breaks if available
\urlstyle{same}
% fallback for those not using the hyperref driver hyperxmp:
\makeatletter
\@ifundefined{xmpquote}{\newcommand{\xmpquote}[1]{#1}}{}
\makeatother
\hypersetup{
  pdftitle={A Counterexample to Arrhenius's Sixth Impossibility Theorem},
  pdfauthor={Claude},
  hidelinks,
  pdfcreator={LaTeX via pandoc}}


\title{A Counterexample to Arrhenius's Sixth Impossibility Theorem}
\author{Claude}
\date{2026}
\begin{document}
\maketitle
\begin{abstract}
Gustaf Arrhenius's sixth impossibility theorem in population ethics, the
one he regards as resting on the weakest and most compelling conditions,
says that no population axiology satisfies Egalitarian Dominance,
General Non-Extreme Priority, Non-Elitism, Weak Non-Sadism, and the Weak
Quality Addition Condition. Teruji Thomas has observed that the
published proof contains a gap, and Elliott Thornley has restated the
last condition with two quantifiers reversed so that the proof goes
through. I show that the gap cannot be closed: the theorem as Arrhenius
states it is false. There is a complete, transitive welfarist axiology,
defined on Arrhenius's own integer-indexed welfare levels, that
satisfies all five of his conditions in their exact formulations. The
axiology ranks populations first by whether their misery exceeds what
their excellent lives can offset, and otherwise by total welfare. It
accepts the Repugnant Conclusion and rejects the Very Repugnant
Conclusion, which is precisely the combination the theorem was designed
to rule out. The repaired theorem is true, and I show what its extra
premise buys: the four remaining conditions force only a
background-relative Very Repugnant Conclusion, and the strengthened
condition is exactly what excludes axiologies on which the quantity of
misery that lives barely worth living can never make up for depends on
how much excellence the world already contains. Every result stated as a
theorem here, including the repaired theorem and Thomas's
reconstructions of the other five, is verified in the Lean 4 proof
assistant; the counterexample was found by a method, described here,
that systematically searches for such axiologies.
\end{abstract}


\setstretch{1.1}
\section{The theorem}\label{sec-intro}

Population ethics has a body of impossibility theorems that do for the
comparison of differently sized populations what Arrow's theorem did for
voting. Each lists a handful of conditions on a betterness ordering of
populations, each condition apparently undeniable on its own, and shows
that no ordering satisfies them all. The problem they formalize is
Parfit's (\citeproc{ref-Parfit1984}{1984}, ch.~17;
\citeproc{ref-Parfit2016}{2016}): that the ways of avoiding the
Repugnant Conclusion each seem to cost something no less repugnant. The
first such theorems were proved by Ng (\citeproc{ref-Ng1989}{1989}) and
Blackorby and Donaldson (\citeproc{ref-BlackorbyDonaldson1991}{1991});
Gustaf Arrhenius has since proved six
(\citeproc{ref-Arrhenius2000a}{Arrhenius 2000b},
\citeproc{ref-Arrhenius2000b}{2000a},
\citeproc{ref-Arrhenius2003}{2003}, \citeproc{ref-Arrhenius2009}{2009},
\citeproc{ref-Arrhenius2011}{2011}), which are now the standard
reference point (\citeproc{ref-Greaves2017}{Greaves 2017}), and the
sixth is the one he regards as best, because its conditions are
logically the weakest. It says:

\begin{quote}
\textbf{The Sixth Impossibility Theorem.} There is no population
axiology which satisfies the Egalitarian Dominance, the General
Non-Extreme Priority, the Non-Elitism, the Weak Non-Sadism, and the Weak
Quality Addition Condition. (\citeproc{ref-Arrhenius2011}{Arrhenius
2011, 9})
\end{quote}

The theorem has been treated as settled. It is stated as a result in
Thornley's work on population prospects
(\citeproc{ref-Thornley2021}{Thornley 2021}), discussed as such by
Carlson (\citeproc{ref-Carlson2022}{2022}), and it is the version of
Arrhenius's results that contributors to the Oxford Handbook of
Population Ethics take as their point of departure
(\citeproc{ref-ArrheniusEtAl2022}{Arrhenius et al. 2022}); Arrhenius's
book, which collects the six theorems, has long circulated in draft
(\citeproc{ref-ArrheniusForthcoming}{Arrhenius, forthcoming}), and I do
not know whether its current version states the condition as the 2011
chapter does. Two readers have noticed that something is wrong with the
proof. Teruji Thomas, in an unpublished reconstruction of all six
theorems, found that the published argument establishes only a weaker
conclusion than the one Arrhenius states, and he noted in print that
``the proof of his favoured sixth theorem contains an error in the
derivation of the `Restricted Quality Addition Condition' (Arrhenius,
2011, Lemma 1.3). As far as I know, one has to slightly strengthen the
premisses of the theorem, in a way unlikely to cause further
controversy'' (\citeproc{ref-Thomas2018}{Thomas 2018, n. 4}; cf.
\citeproc{ref-Thomas2016}{Thomas 2016}). Elliott Thornley, following
Thomas, restates the Weak Quality Addition Condition with its first two
quantifiers reversed, remarks that ``the Sixth Impossibility Theorem
actually requires the slightly stronger condition stated here,'' and
adds that ``the stronger version remains a compelling adequacy
condition'' (\citeproc{ref-Thornley2021}{Thornley 2021, n. 7}).

Thomas's remark already implies that the theorem, and not only its
proof, has to change; but neither he nor Thornley shows that it does,
and the field has gone on citing the original statement with the
strengthening treated as a formality. I will show that the strengthening
is not a formality. The theorem as Arrhenius states it is false. There
is a complete, transitive, anonymous welfarist axiology, defined on
exactly the integer-indexed welfare levels that Arrhenius's framework
assumes, that satisfies all five of his conditions in the exact
formulations he gives them. I call it the \emph{capacity view}. It ranks
populations first by whether the misery they contain exceeds what their
excellent lives can offset, and, when it does not, by total welfare. It
accepts the Repugnant Conclusion, as total utilitarianism does, and it
rejects the Very Repugnant Conclusion, which is the combination of
attitudes that the move from Arrhenius's earlier theorems to the sixth
was designed to make untenable. So the sixth theorem does not refute the
position it was built to refute.

The strengthened condition that Thomas and Thornley supply does make the
theorem true; Thomas proved the strengthened theorem in his note, and I
have checked his proof. What I want to add is an account of what the
strengthening buys, since the capacity view makes the difference between
the two conditions concrete. In its original form the Weak Quality
Addition Condition says that, for any population, there is some quantity
of very bad lives such that adding those lives together with any number
of lives barely worth living is no better than adding some number of
excellent lives. The strengthened condition fixes the quantity of very
bad lives first and then quantifies over populations. The difference is
whether the amount of misery that lives barely worth living can never
make up for is allowed to depend on how much excellence the population
already contains. The capacity view says that it is. The four other
conditions, I will show, cannot force it not to be: what they entail is
a Very Repugnant Conclusion relative to a background that may be chosen
after the quantity of misery is fixed, and nothing stronger. The
repaired theorem is therefore a theorem about background-relative
repugnance, and whether its new premise is as ``compelling'' as the old
one is a question about whether capacity views deserve to be ruled out
by an adequacy condition rather than by argument.

The paper also has a methodological point. The counterexample was not
found by inspection. It was found by looking for a quantity that none of
the four non-conclusion conditions can decrease, and then building an
axiology that gives that quantity lexical priority.
Section~\ref{sec-method} describes the method, which applies to any
impossibility theorem whose conditions are ``transfer'' conditions of
the kind Arrhenius uses, and which turns the search for countermodels
into a routine check. Everything in the paper has been verified in the
Lean 4 proof assistant: that the capacity view satisfies the five
conditions, that it violates the strengthened condition, and that
Thomas's reconstructions of Theorems 1 through 5 and of the corrected
sixth theorem are valid. Section~\ref{sec-lean} describes the
verification, and the source is available with the paper.

\section{The framework and the five conditions}\label{sec-conditions}

Arrhenius's framework is this (\citeproc{ref-Arrhenius2011}{Arrhenius
2011, 4--7}). A life has a welfare level; welfare levels are
quasi-ordered; and there is a neutral level, with lives above it worth
living and lives below it worth not living. Arrhenius assumes
\emph{Discreteness}: between any two welfare levels there are only
finitely many others. (He argues that nothing turns on this, since if
welfare is dense one can pass to a discrete subset of levels, as
fine-grained as one likes, on which the conditions remain plausible
(\citeproc{ref-Arrhenius2011}{Arrhenius 2011, 6}). I will return to the
point in Section~\ref{sec-repair}.) Given Discreteness, welfare levels
can be indexed by integers, with 0 the neutral level, so that the level
\(w+1\) is ``slightly higher'' than \(w\). A \emph{population} is a
finite set of lives. Following Thomas (\citeproc{ref-Thomas2016}{2016})
I will identify a population with its welfare distribution, the function
that assigns to each level the number of lives at that level, and write
\(n \cdot x\) for a population of \(n\) lives all at level \(x\), and
\(P + Q\) for the union of two populations.\footnote{Arrhenius's
  populations are sets of particular lives, and his conditions are
  stated with the rider ``other things being equal,'' which abstracts
  from non-welfare differences between them. Since every condition is
  invariant under replacing lives by others at the same level, nothing
  is lost by working with distributions. Arrhenius's own exact
  formulations use the notation \(A \subset W_x\), \(N(A) = n\) for what
  I write \(n \cdot x\).} An \emph{axiology} is a relation ``at least as
good as'' on populations, assumed to be reflexive and transitive but not
necessarily complete; ``better than'' is the strict part.

Three further parameters appear in the conditions. There is a ``very
high'' welfare level, which I write \(A\); a ``very low positive'' range
of levels, which Arrhenius takes to be \(R(1,y)\), the levels from 1 up
to some \(y\); and a ``very negative'' level. In the exact formulations
these are not fixed in advance but are quantified over inside each
condition, which is one of the things that makes the formulations exact;
a condition that says ``there is a very high level \(u\) such
that\ldots{}'' does not presuppose that we know where ``very high''
begins. The reader who wants the informal versions will find them
alongside the exact ones in Arrhenius
(\citeproc{ref-Arrhenius2011}{2011, 7--9}); I give the exact ones, in
distribution notation, since the result depends on them. All welfare
levels are integers, all population sizes are non-negative integers, and
ranges \(R(x,y)\) consist of at least three levels.

\begin{quote}
\textbf{Egalitarian Dominance (ED).} If \(N(A) = N(B)\), every member of
\(B\) has welfare below \(x\), and \(A = n \cdot x\), then \(A\) is
better than \(B\).
\end{quote}

\begin{quote}
\textbf{General Non-Extreme Priority (GNEP).} For any level \(z\) there
are a positive level \(u\), a positive range \(R(1,y)\) with \(u > y\),
and a number \(n > 0\) such that for any \(x \geq u\), any
\(B \subset R(1,y)\) with \(N(B) = n\), and any population \(E\):
\(\;n \cdot x + 1\cdot z + E\) is at least as good as
\(B + 1 \cdot (z+1) + E\).
\end{quote}

\begin{quote}
\textbf{Non-Elitism (NE).} For any levels \(x, y\) with \(x - 1 > y\)
there is \(n > 0\) such that for any \(D \subset R(y, x)\):
\(\;(n+1) \cdot (x-1) + D\) is at least as good as
\(1 \cdot x + n \cdot y + D\).
\end{quote}

\begin{quote}
\textbf{Weak Non-Sadism (WNS).} There are a level \(x < 0\) and a number
\(n\) such that for any level \(y > 0\), any \(B \subset W_y\), and any
population \(C\): \(\;B + C\) is at least as good as \(n \cdot x + C\).
\end{quote}

\begin{quote}
\textbf{Weak Quality Addition (WQA).} For any population \(X\) there are
a level \(x < 0\), positive ranges \(R(u,v)\) and \(R(1,y)\) with
\(u > y\), and numbers \(n, m > 0\) such that for any \(z \geq u\) and
any \(B \subset R(1,y)\): \(\;n \cdot z + X\) is at least as good as
\(B + m \cdot x + X\).
\end{quote}

The informal content is familiar. ED says that a perfectly equal
population is better than any same-sized population every member of
which is worse off than its members. GNEP says that lowering one
person's welfare by a step is never categorically worse than failing to
raise a fixed number of others from very low positive to very high
welfare, whatever else is true of the population. NE says that one
person's small loss can always be compensated by a sufficient number of
others' gains from a lower level to the same level she falls to, at
least when the rest of the population lies between those levels. WNS
says that there is some number of lives at some very negative level
whose addition is no better than the addition of any number of lives
with positive welfare. And WQA says that for any population there is
some quantity of very bad lives such that adding them, together with any
number of lives barely worth living, is no better than adding some
number of very good lives.

Three features of WQA matter below. First, the conclusion compares
additions to a \emph{background} \(X\), and the witnesses (\(x\), \(m\),
\(u\), \(y\), \(n\)) are chosen after \(X\) is given: the condition
begins ``for any population \(X\) there are.'' Second, \(B\) ranges over
all populations in the low range, of every size, so that the condition
says the \(m\) bad lives cannot be compensated by any number of marginal
lives. Third, the condition is a weakening of Arrhenius's Quality
Addition Condition, which drops the \(m\) bad lives and says that for
any \(X\) some number of very good lives added to \(X\) is at least as
good as any number of marginal lives added to \(X\). Quality Addition is
the denial of what Thomas calls Repugnant Addition, a
background-relative form of the Repugnant Conclusion, and total
utilitarians reject it. Weak Quality Addition is the denial of a
background-relative form of the Very Repugnant Conclusion, that for any
population of excellent lives there is a better population consisting of
some very bad lives together with enough lives barely worth living.
Arrhenius weakened the condition because the Very Repugnant Conclusion
is, as he and his coauthors put it, ``intuitively more compelling'' as a
thing to avoid than the Repugnant Conclusion
(\citeproc{ref-ArrheniusBudolfsonSpears2021}{Arrhenius, Budolfson, and
Spears 2021, 119}): it is meant to bind even those who, like Tännsjö
(\citeproc{ref-Tannsjo2002}{2002}) and Huemer
(\citeproc{ref-Huemer2008}{2008}), have made their peace with the
Repugnant Conclusion, a group that has lately grown
(\citeproc{ref-ZuberEtAl2021}{Zuber et al. 2021}) as it has become clear
how little is needed to derive the conclusion
(\citeproc{ref-SpearsBudolfson2021}{Spears and Budolfson 2021}).

\section{The capacity view}\label{sec-capacity}

Fix a ``very high'' level \(A \geq 4\), so that the very low positive
levels 1, 2, 3 lie strictly below it. For a population \(P\) let

\begin{itemize}
\tightlist
\item
  \(H(P)\) be the number of lives in \(P\) at or above level \(A\) (the
  \emph{excellent} lives);
\item
  \(T^-(P)\) be the total welfare of the lives in \(P\) with negative
  welfare, a non-positive number;
\item
  \(T(P)\) be the total welfare of \(P\).
\end{itemize}

Let \(I(P) = H(P) + T^-(P)\), and for a fixed \emph{baseline capacity}
\(c \geq 0\) let \(J(P) = \min(0, I(P) + c)\). The quantity \(I\)
measures the population's misery against its excellence, counting each
excellent life as able to offset one unit of negative welfare; \(c\) is
an amount of misery that any population, including the empty one, can
absorb. \(J\) is zero when the excellent lives, together with the
baseline, cover the misery, and negative, equal to the uncovered misery,
when they do not. The \emph{capacity view} (with parameters \(A\) and
\(c\)) ranks populations lexically by \((J, T)\):

\begin{quote}
\(P\) is at least as good as \(Q\) if and only if either
\(J(P) > J(Q)\), or \(J(P) = J(Q)\) and \(T(P) \geq T(Q)\).
\end{quote}

This is a complete and transitive ordering, since it is the
lexicographic order on pairs of integers. It is welfarist and anonymous,
since it depends only on the distribution. In words: a population whose
misery exceeds its capacity to redeem it is worse than any population
whose misery does not, and among populations whose misery exceeds their
capacity the one with the smaller excess is better; among populations
within capacity, total welfare decides. The exchange rate of one
excellent life per unit of misery is a choice of convenience; any
positive weight on \(H\) works, with the witness \(n\) in GNEP and NE
adjusted. The simplest member of the family has \(c = 0\), and the
reader may take that as the official countermodel; the baseline matters
only for the discussion in Section~\ref{sec-repair}.

The view has some features its critics will want to know about at once.
It is lexical, and so discontinuous: a single unit of uncovered misery
outweighs any quantity of total welfare. With \(c = 0\) this has a
consequence worth stating in its starkest form: a population of a
billion lives at level 3 together with one life at level \(-1\) has
\(J = -1\), and is therefore worse than the empty population. Lives
below the very high level cannot redeem misery at all, however many
there are. With \(c \geq 1\) the consequence disappears, since the
baseline absorbs the one unit of misery, and I return to what this shows
in Section~\ref{sec-repair}. It implies the Repugnant Conclusion, since
populations containing no misery are ranked by total welfare, so that
for any \(n\) lives at level \(A\) there is a larger population of lives
at level 3 that is better. It rejects the Very Repugnant Conclusion: for
\(m|z| > c\), the population \(m \cdot z + N \cdot 3\) has
\(I + c = mz + c < 0\) and so is worse than \(n \cdot A\), which has
\(J = 0\), for every \(N\). And it is not meant as a proposal about how
to rank populations. It is a proof that the five conditions can be
jointly satisfied, and its interest lies in what it reveals about them.

\textbf{Theorem 1.} For every \(A \geq 4\) and \(c \geq 0\), the
capacity view satisfies ED, GNEP, NE, WNS and WQA in the exact
formulations of Section~\ref{sec-conditions}.

I give the arguments at the level of detail a reader needs to see why
each condition holds; the fully formal versions are the Lean proofs
described in Section~\ref{sec-lean}. The key observation is that every
one of the first four conditions compares two populations that differ by
a \emph{transfer}: some lives are moved between levels, or some are
added. Write \(\Delta I\) and \(\Delta T\) for the change in \(I\) and
\(T\) from the population the condition says is no better to the one it
says is at least as good. If \(\Delta I \geq 0\) and \(\Delta T \geq 0\)
the capacity view agrees with the condition, since \(J\) is monotone in
\(I\) (the baseline \(c\) cancels from every such comparison); and if
\(\Delta I \geq 0\) and \(\Delta T > 0\) it agrees with the strict
version. Background populations drop out, because \(I\) and \(T\) are
additive over union.

\emph{Egalitarian Dominance.} \(A = n \cdot x\) and every member of
\(B\) is below \(x\), with \(N(B) = n\). Then \(T(A) = nx > T(B)\). For
\(I\): if \(x \geq A\) then \(H(A) = n \geq H(B)\), and if \(x < A\)
then \(H(B) = 0 = H(A)\); and \(T^-(A) \geq T^-(B)\), since if \(x < 0\)
every member of \(B\) is negative and below \(x\), so
\(T^-(B) = T(B) < nx = T^-(A)\), while if \(x \geq 0\) then
\(T^-(A) = 0 \geq T^-(B)\). So \(\Delta I \geq 0\) and \(\Delta T > 0\).

\emph{General Non-Extreme Priority.} Given \(z\), choose \(u = A\),
\(y = 3\), \(n = 2\). The condition compares \(B + 1 \cdot (z+1) + E\)
with \(2 \cdot x + 1 \cdot z + E\), where \(x \geq A\) and \(B\) is two
lives in \(\{1,2,3\}\). The two excellent lives raise \(H\) by 2; the
person who moves from \(z+1\) to \(z\) lowers \(H\) by 1 if
\(z + 1 = A\), and lowers \(T^-\) by 1 if \(z + 1 \leq 0\); \(B\)
contributes nothing to \(H\) or \(T^-\). So
\(\Delta I \geq 2 - 1 - 1 = 0\). And
\(\Delta T = 2x + z - T(B) - (z+1) \geq 2A - 6 - 1 > 0\).

\emph{Non-Elitism.} Given \(x - 1 > y\), choose \(n = 1\). The condition
compares \(1\cdot x + 1 \cdot y + D\) with \(2 \cdot (x-1) + D\), for
\(D \subset R(y,x)\). \(\Delta T = x - 2 - y \geq 0\). For \(I\) there
are cases. If \(x - 1 \geq A\) the two lives at \(x-1\) are excellent
and \(\Delta H \geq 0\), while \(\Delta T^- \geq 0\) since only the life
at \(y\) can be negative and it is raised. If \(x < A\) then
\(\Delta H = 0\) and a short calculation shows \(\Delta T^- \geq 0\):
the only case needing care is \(x = 0\), where two lives at \(-1\)
replace one at \(0\) and one at \(y \leq -2\), and
\(\Delta T^- = -2 + |y| \geq 0\). The remaining case is \(x = A\), where
one excellent life is lost: \(\Delta H = -1\). If \(y < 0\) then the
life raised from \(y\) gives \(\Delta T^- = |y| \geq 1\), so
\(\Delta I \geq 0\). If \(y \geq 0\) then \(D \subset R(y, A)\) contains
no negative lives, so both populations have \(T^- = 0\), both have
\(I + c \geq 0\), both have \(J = 0\), and the comparison is decided by
\(T\), which does not decrease. This is the one place where \(I\) can
fall under a licensed transfer, and it can fall only in a population
whose \(J\) is already zero and stays zero. The restriction on \(D\) in
Non-Elitism is what makes this work, and it is the restriction that
Arrhenius's General Non-Elitism Condition, used in his fifth theorem,
drops; the capacity view violates General Non-Elitism, as it must, since
the fifth theorem is valid.

\emph{Weak Non-Sadism.} Choose \(x = -1\) and \(n = 1\). The condition
compares \(1 \cdot (-1) + C\) with \(B + C\), where \(B\) is any number
of lives at some \(y > 0\). \(\Delta I = H(B) + 1 \geq 1\) and
\(\Delta T = T(B) + 1 \geq 1\).

\emph{Weak Quality Addition.} Given \(X\), choose \(x = -1\), \(u = A\),
\(y = 3\), \(n = 1\), and \(m\) large enough that \(m > I(X) + c\);
concretely \(m = \max(0, I(X) + c) + 1\). For any \(z \geq A\) and any
\(B \subset R(1,3)\): \(I(1 \cdot z + X) = 1 + I(X)\), while
\(I(B + m \cdot (-1) + X) + c = I(X) + c - m < 0\), since \(B\)
contributes nothing to \(I\). So
\(J(B + m\cdot(-1) + X) = I(X) + c - m < \min(0, 1 + I(X) + c) = J(1 \cdot z + X)\),
and \(1 \cdot z + X\) is better than \(B + m \cdot (-1) + X\) for every
\(B\), whatever its size.

The last argument displays the mechanism. Weak Quality Addition lets the
quantity of misery depend on the background, and the capacity view uses
exactly that freedom: it chooses enough misery to exceed the
background's capacity, after which no quantity of marginal lives can
help, because marginal lives contribute to \(T\) and not to \(I\). The
four other conditions never require the view to compare a population
that has exceeded its capacity favourably with one that has not,
because, as the case analysis shows, none of the transfers they license
lowers \(I\) except in populations with no misery at all.

Since the capacity view satisfies the five conditions, the Sixth
Impossibility Theorem is false. Arrhenius's argument, I should stress,
is not invalid at some obscure step; it is invalid at exactly the step
Thomas identified. In the proof of his Lemma 1.3, which derives a
Restricted Quality Addition Condition from Weak Quality Addition and an
intermediate condition, step (12) applies Weak Quality Addition to the
background population \(A_1 \cup X\) using witnesses that the condition
supplies for the background \(X\)
(\citeproc{ref-Arrhenius2011}{Arrhenius 2011, 18--19}). The population
\(A_1\) consists of excellent lives whose number depends, through the
intermediate condition, on those very witnesses. On the capacity view
the step fails for a simple reason: adding \(A_1\) raises the
background's capacity, and the quantity \(m\) of misery that sufficed to
exceed the capacity of \(X\) no longer exceeds the capacity of
\(A_1 \cup X\), so that in the enlarged background the marginal lives
win.

\section{The repair and what it costs}\label{sec-repair}

Thomas states the conclusion that the four conditions do yield, and
Thornley states the condition whose denial it is. In Thomas's notation
(\citeproc{ref-Thomas2016}{2016, 11}), the conclusion Arrhenius intended
is

\begin{quote}
\textbf{Very Repugnant Addition (VRA).} For some population \(P\), any
level \(z < 0\), and any \(m, n \geq 1\), there is \(N\) such that
\(P + m \cdot z + N \cdot 3\) is better than \(P + n \cdot A\),
\end{quote}

and the conclusion the proof delivers is

\begin{quote}
\textbf{Very Repugnant Addition* (VRA*).} For any level \(z < 0\) and
any \(m \geq 1\), there is a population \(P\) such that for any
\(n \geq 1\) there is \(N\) with \(P + m \cdot z + N \cdot 3\) better
than \(P + n \cdot A\).
\end{quote}

The difference is that in VRA* the background \(P\) may depend on \(z\)
and \(m\). (In Thomas's fixed-threshold vocabulary, and for a complete
ordering, VRA* is equivalent to the denial of the condition WQA* below;
for an incomplete ordering the denial of WQA* is weaker, since WQA* asks
for positive comparisons.) Thomas writes that VRA* ``is slightly weaker
than VRA, because \(P\) is allowed to depend on \(z\) and \(m\). I think
that VRA* is not much less `repugnant' than VRA, but then again it is
not easy to understand the difference between them''
(\citeproc{ref-Thomas2016}{Thomas 2016, 11}). Thornley's strengthened
condition, which entails the denial of VRA*, is:

\begin{quote}
\textbf{Weak Quality Addition* (WQA*).} There are a level \(x < 0\) and
a number \(m > 0\) such that for any population \(X\) there are a
positive level \(u\), a range \(R(1,y)\) with \(u > y\), and \(n > 0\)
such that for any \(z \geq u\) and any \(B \subset R(1,y)\):
\(\;n \cdot z + X\) is at least as good as \(B + m \cdot x + X\).
\end{quote}

Thomas's corrected theorem, in his reconstruction, is that ED, GNEP, NE
and WNS entail VRA*, from which it follows that ED, GNEP, NE, WNS and
WQA* are jointly unsatisfiable. I have verified his proof, and also
proved the unsatisfiability directly for the conditions in Arrhenius's
exact formulations rather than Thomas's fixed-threshold ones
(Section~\ref{sec-lean}), so the repaired theorem is true in both forms.
The capacity view violates WQA*: given \(x\) and \(m\), take \(X\) to be
\(m|x|\) lives at level \(A\), so that the \(m\) lives at \(x\) are
exactly covered and \(I(B + m \cdot x + X) = 0\); then both sides have
\(J = 0\), total welfare decides, and enough lives at level 1 in \(B\)
make the right-hand side better. (The baseline only helps: with
\(c > 0\) the misery is covered with room to spare.)

The question Thomas raises, whether VRA* is much less repugnant than
VRA, is the question whether WQA* is as compelling an adequacy condition
as WQA, and the capacity view lets us say exactly what is at stake. The
two conditions agree on every axiology for which the following holds:

\begin{quote}
\textbf{Compensation.} For any population \(Y\), any level \(z < 0\) and
any \(m\), there is \(K\) such that \(Y + m \cdot z + K \cdot 3\) is at
least as good as \(Y\).
\end{quote}

Compensation says that any finite quantity of misery added to any
population can be made up for by adding enough lives barely worth
living. Total utilitarianism satisfies it, as do critical-level views
with a critical level below 3
(\citeproc{ref-BlackorbyBossertDonaldson2005}{Blackorby, Bossert, and
Donaldson 2005}; \citeproc{ref-Broome2004}{Broome 2004}), prioritarian
views, and in general any view that is Archimedean in the trade between
misery and marginal lives. For such views WQA* and WQA collapse into
each other, and both collapse into Quality Addition:

\textbf{Proposition 2.} If an axiology satisfies Compensation and WQA,
it satisfies Quality Addition in Arrhenius's exact form: for any \(X\)
there are a positive level \(u\), a range \(R(1,y)\) with \(u > y\), and
\(n > 0\) such that \(n \cdot z + X\) is at least as good as \(B + X\)
for every \(z \geq u\) and every \(B \subset R(1,y)\).

\emph{Proof.} Given \(X\), let \(x, u, y, n, m\) be the witnesses WQA
supplies. Take any \(z \geq u\) and any \(B \subset R(1,y)\).
Compensation applied to the population \(X + B\) gives \(K\) with
\(X + B + m \cdot x + K \cdot 3\) at least as good as \(X + B\). Since
\(y \geq 3\), the population \(B + K \cdot 3\) lies in \(R(1,y)\), so
WQA gives \(n \cdot z + X\) at least as good as
\(B + K \cdot 3 + m \cdot x + X\). Transitivity gives \(n \cdot z + X\)
at least as good as \(B + X\). \(\square\)

The same argument with WQA* in place of WQA goes through unchanged. So
for any axiology satisfying Compensation, the weakening of Quality
Addition to Weak Quality Addition that distinguishes the sixth theorem
from the fourth buys nothing: a Compensation-satisfying axiology that
accepts Weak Quality Addition, in either form, accepts Quality Addition.
(The converse needs only that adding very bad lives never improves a
population, which every view in the vicinity grants.) Put together with
Thomas's fourth theorem, which says that ED, GNEP, NE and WNS entail
Repugnant Addition, the denial of Quality Addition, this gives a
corollary worth stating on its own, in Thomas's fixed-threshold setting
where the two theorems share their vocabulary:

\textbf{Corollary 3.} ED, GNEP, NE, WNS, WQA and Compensation are
jointly unsatisfiable. That is: given the four conditions that the sixth
theorem shares with the fourth, Weak Quality Addition entails that some
finite quantity of misery cannot be made up for by any number of lives
barely worth living.

The weakening of Quality Addition to Weak Quality Addition therefore has
content only for axiologies that violate Compensation, and, given the
other four conditions, it \emph{commits} one to violating Compensation.
Those are the axiologies at which the sixth theorem is aimed. They
include lexical views on which the badness of a very bad life is of a
different order from the goodness of a marginal one, such as the Lexical
Totalism that Thomas (\citeproc{ref-Thomas2018}{2018}) and Thornley
(\citeproc{ref-Thornley2021}{2021}) discuss and that Thornley
(\citeproc{ref-Thornley2022}{2022}) elsewhere argues faces a dilemma of
its own, and they include the capacity view.

Among Compensation-violating axiologies, WQA and WQA* come apart, and
the capacity view is the shape of the gap. WQA* says that the quantity
of misery that marginal lives cannot make up for can be fixed once,
independently of the population to which the misery is added. WQA says
only that for each population there is such a quantity. The capacity
view takes the second option: the quantity depends on how much
excellence the population already contains, because the view treats
misery as redeemable by excellent lives, one for one, and as
irredeemable by marginal lives. In a population rich in excellent lives,
any fixed quantity of misery is covered, and once covered it is subject
to the ordinary totalist trade-off against marginal lives; in a
population poor in excellent lives, the same quantity is uncovered, and
no quantity of marginal lives touches it.

Is that structure one an adequacy condition should rule out? I do not
think the question has an obvious answer, and I want to resist two quick
ones.

The first quick answer is that the capacity view is obviously absurd, so
that WQA* is as harmless as Thornley says. The view has features nobody
would defend as stated: the exchange rate between an excellent life and
a unit of misery is arbitrary, the lexical threshold is sharp, and the
view was constructed to satisfy five conditions rather than to be true.
With \(c = 0\) it also says that a billion lives barely worth living
together with one slightly bad life is worse than no lives at all. That
last consequence is worth separating from the others, because it is not
a commitment of the capacity structure. The baseline parameter removes
it: with \(c \geq 1\) the empty population absorbs a unit of misery and
the billion marginal lives then count in its favour, while every one of
the five conditions continues to hold and the Very Repugnant Conclusion
continues to fail for quantities of misery exceeding the baseline. Nor
is the consequence peculiar to the \(c = 0\) view among the sixth
theorem's escapes. By Corollary 3, anyone who accepts the five
conditions denies Compensation: there is some quantity of misery that no
number of lives barely worth living can make up for. The \(c = 0\) view
denies it at every background, including the empty one. So does Lexical
Totalism in the form Thornley gives it, on which ``adding a combination
of very negative welfare lives and barely positive welfare lives means
adding \((r, s)\) with \(r < 0\)'' to a value vector whose first
coordinate dominates lexically (\citeproc{ref-Thornley2021}{Thornley
2021, sec. 4}): a very bad life subtracts a higher good that no number
of barely-worth-living lives, which add only lower goods, can restore,
so there too misery plus marginal lives is worse than nothing. The five
conditions do not force this, as the \(c \geq 1\) view shows; but the
one escape from the sixth theorem in the published literature has it,
and the capacity view with a positive baseline is, as far as I know, the
first axiology satisfying the conditions that does not.

A second feature invites the same treatment. In the view as defined,
redemptive power belongs only to lives at or above the very high level
\(A\); lives between the very low range and \(A\) have none, so that for
large \(A\) a billion lives at level \(A - 1\) cannot cover a single
unit of misery beyond the baseline, while for \(A = 4\) ``excellence''
begins one step above lives barely worth living. Neither horn is forced
either. The step at \(A\) is an artefact of taking the simplest
potential. A graded version, in which the weight of a life rises
linearly from 0 at level 3 to 1 at level \(A\) and the Non-Elitism
witness is taken to be \(A - 3\), satisfies all five conditions as well;
I have checked this numerically for \(A\) up to 10, though the
machine-checked proof is for the step version. On the graded view every
life above the barely-worth-living range contributes to redeeming
misery, in proportion to how far above it lies, and only the
barely-worth-living lives contribute nothing. That is the version of the
view that the gloss ``excellence redeems, numbers do not'' properly
describes, and in that form the structural idea is not absurd. It is
close to the intuition that the Very Repugnant Conclusion is worse than
the Repugnant Conclusion: someone who accepts that enough marginal lives
can outweigh a population of excellent lives, but denies that they can
outweigh it when a quantity of misery has been thrown in, is treating
misery as a dimension of value that marginal lives do not reach. The
capacity view is the simplest way of making that attitude consistent
with the four other conditions, and the fact that it must let the
uncompensable quantity of misery depend on the background is a discovery
about the attitude, not an artefact of the construction.

The second quick answer is that VRA* is as repugnant as VRA, so the
weakening does not matter. Here the capacity view gives a reason to
disagree with Thomas's first instinct, though I should be clear about
whose reason it is. VRA* says that for any given quantity of misery
there is a background against which adding that misery together with
enough marginal lives beats adding excellent lives. On the capacity view
the background that does the work is one rich enough in excellent lives
to cover the misery, and what VRA* then asserts is that, in a population
whose excellent lives already redeem its misery, enough additional
marginal lives beat additional excellent lives. That is the Repugnant
Conclusion again, stated for a population with some covered misery in
it, and someone who has accepted the Repugnant Conclusion has no new
reason to find it repugnant. That is how the capacity view describes its
own verdict. A critic of the Very Repugnant Conclusion who is not a
capacity theorist will describe the same verdict differently: in some
sufficiently good world, adding lives of torment together with enough
lives barely worth living is better than adding excellent lives, and the
goodness of the surrounding world does not make the lives of torment
less tormented. I do not think the capacity view has an answer to that
critic that does not presuppose the view. What it has is a diagnosis of
what the critic is committed to. The critic's complaint is against VRA*,
which the four conditions entail; so the critic must reject one of ED,
GNEP, NE and WNS, and Weak Quality Addition in any form will not help,
because the critic's complaint is about the background-relative
conclusion that the four conditions reach without it. The sixth theorem,
repaired, tells the critic of the Very Repugnant Conclusion that the
fight is elsewhere. The unrepaired theorem told the critic, wrongly,
that there was no consistent position to fight from.

So the corrected sixth theorem is a theorem about background-relative
repugnance. What ED, GNEP, NE and WNS jointly force is that the
Repugnant Conclusion holds relative to any background that has redeemed
its misery. They do not force the Very Repugnant Conclusion, and they do
not force it relative to a fixed background. Anyone who wants the
stronger result must add WQA* as a premise, and the only thing WQA* adds
to WQA is the stipulation that the quantity of misery which marginal
lives cannot make up for be the same for every background. Proposition 2
shows that this stipulation has no work to do against Archimedean views,
which accept or reject both conditions together with Quality Addition;
its content is the exclusion of views on which the uncompensable
quantity of misery depends on the background, of which the capacity view
is the simplest instance.

Two remarks on scope. First, the result does not depend on the denial of
Finite Fine-Grainedness. The known counterexamples to the sixth theorem,
Lexical Totalism among them, work by giving welfare levels a lexical
structure, so that no finite sequence of slight steps leads from a very
high level to a very low one; Thornley says plainly that ``the theorem
is only true given Finite Fine-Grainedness''
(\citeproc{ref-Thornley2021}{Thornley 2021, sec. 3}), and he is right
about the strengthened theorem. The capacity view lives on Arrhenius's
integer levels, where every level is finitely many slight steps from
every other, and it satisfies the original conditions there. The lexical
structure in the view is in the ranking of populations, not in the
ranking of lives, and the four conditions other than WQA, which are the
ones that use fine-grainedness, do not reach it. Second, the capacity
view also satisfies several conditions not in the theorem: Dominance
Addition, Non-Sadism, and Mere Addition among them. I note this only to
forestall the thought that the view satisfies the five conditions by
being pathological elsewhere. Where it is pathological is where it was
designed to be, in the lexical priority of uncovered misery.

\section{How the counterexample was found}\label{sec-method}

The capacity view was not found by trying candidate axiologies against
the conditions. It was found by a search that I think is worth
describing, because it applies to any impossibility theorem of this
type, because it explains why the sixth theorem fails where the others
do not, and because it is cheap enough that it should be run on any new
theorem of the kind.

Arrhenius's conditions, other than the Quality conditions, are
\emph{transfer conditions}. Each says that a population obtained from
another by a specified change, moving some lives between levels or
adding some lives, is at least as good as the original, possibly
strictly, possibly only when the background satisfies a restriction.
Read as licences, they say which transfers an axiology must count as
non-worsening. An impossibility proof is a chain of licensed transfers
from one population to another, with a licensed strict transfer
somewhere in the chain, ending at a population that the conclusion's
denial says is no better than the start. The conclusion's denial is the
odd condition out: it is not a transfer licence but a demand that
certain comparisons go a certain way, and an axiology satisfies it by
\emph{refusing} a comparison that the chain would otherwise force.

This suggests looking for a quantity \(\Phi\) on populations that no
licensed transfer lowers, and building the axiology that ranks
populations first by \(\min(0, \Phi)\) and then by total welfare. If
\(\Phi\) is additive over union, with a weight \(w_\ell\) for each
welfare level \(\ell\), then backgrounds drop out and each transfer
condition becomes a linear inequality on the weights: General
Non-Extreme Priority with witness \(G\) requires
\(w_{z-1} + G w_x \geq w_z + G w_y\), Non-Elitism with witness \(n\)
requires \((n+1) w_{x-1} \geq w_x + n w_y\), Weak Non-Sadism with
witness \(D\) requires \(m w_x \geq D w_Z\), and Egalitarian Dominance
requires the weights to be non-decreasing in the level. Two further
constraints come from the lexical layer. One transfer lowers every
candidate \(\Phi\) of this kind: Non-Elitism at the very high level,
when the lives being raised have non-negative welfare, removes an
excellent life and adds nothing in the negative range. That transfer is
licensed only in backgrounds restricted to levels at or above the raised
lives', hence to populations with no misery, so it can be neutralized by
capping the lexical layer at zero, provided every population without
negative lives has \(\Phi \geq 0\), which means \(w_\ell \geq 0\) for
\(\ell \geq 0\). That is why the layer is \(\min(0, \Phi)\) rather than
\(\Phi\). And refusing a conclusion of the form ``for all \(N\),
\(m \cdot z + N \cdot 3 + P\) is better than \(n \cdot A + P\)''
requires \(\Phi(m \cdot z + N \cdot 3) < 0\) for all \(N\), hence
\(w_3 \leq 0\) and so, with monotonicity and the cap,
\(w_0 = w_1 = w_2 = w_3 = 0\), together with \(w_z < 0\).

For fixed witness sizes this is a linear feasibility problem in a dozen
variables, and a general-purpose solver settles it instantly; the script
that sets it up and runs it, for all six theorems, accompanies the
paper. For the conditions of the sixth theorem, with every witness set
to 1, it returns \(w_\ell = \ell\) for \(\ell < 0\), \(w_\ell = 0\) for
\(0 \leq \ell < A\), and \(w_\ell = 1\) for \(\ell \geq A\): that is,
\(\Phi = H + T^-\), the capacity view. For the conditions of the fifth
theorem, which replace Non-Elitism by General Non-Elitism and Weak
Non-Sadism by Dominance Addition, the system is infeasible for every
small witness size, because the unrestricted background lets the
transfer at the very high level occur in populations with uncovered
misery, where it must not lower \(\Phi\), and that forces
\(w_A \leq 0\), which the remaining conditions then propagate into a
contradiction with the refusal. For the first four theorems the method
is uninformative, and the reason is worth stating correctly. Their
conclusions involve no misery: they say that enough lives at level 3
beat any number at level \(A\), against a fixed background. Since total
welfare grows without bound as lives at level 3 are added, an axiology
of the shape under consideration can refuse such a conclusion only
through the lexical layer, which requires \(\Phi(N \cdot 3)\) to go to
\(-\infty\), that is, \(w_3 < 0\); and the cap requires \(w_3 \geq 0\).
So no axiology of this shape refuses a misery-free conclusion, whatever
the transfer conditions, and the system is infeasible before the
conditions are consulted. The sixth theorem is the only one of the six
whose conclusion involves misery, and the fifth is the only other one
that could in principle be refused this way; there the transfer
conditions do the excluding. The search finds the gap in the sixth
theorem because the gap is there, and finds nothing in the fifth because
General Non-Elitism closes it (the system with General Non-Elitism
removed and Non-Elitism restored is feasible again).

The method has limits. It finds countermodels of one shape, lexical in
an additive potential and totalist below it, and an impossibility
theorem can fail for other reasons; infeasibility of the system for
small witnesses is evidence rather than proof that no countermodel of
this shape exists, since the witnesses are existentially quantified and
larger ones might help. What the method does deliver, when it succeeds,
is a candidate whose verification is routine; when it fails on a
conclusion of the right shape, as for the fifth theorem, it shows which
condition does the excluding. In this case it delivered the capacity
view in the first run, and the verification took longer than the search.

\section{Machine verification}\label{sec-lean}

Impossibility theorems in population ethics are proved by long chains of
inequalities with many quantifier dependencies, and the error in the
sixth theorem is a quantifier-dependence error: a witness chosen for one
background is used for another. This is the kind of error a proof
assistant catches, and the kind that referees do not. Social choice
theory learned this some time ago: Arrow's theorem and its relatives
have been formally verified (\citeproc{ref-Nipkow2009}{Nipkow 2009}),
and SAT solvers have been used both to re-prove known impossibility
theorems (\citeproc{ref-TangLin2009}{Tang and Lin 2009}) and to discover
new ones (\citeproc{ref-GeistEndriss2011}{Geist and Endriss 2011};
\citeproc{ref-BrandtGeist2016}{Brandt and Geist 2016}). Population
axiology has, so far as I can find, one mechanized result, a
formalization of Parfit's Mere Addition Paradox in a deontic logic
(\citeproc{ref-ParentBenzmuller2024}{Parent and Benzmüller 2024}), and
none of Arrhenius's theorems had been checked. I have therefore
formalized the results of this paper in Lean 4
(\citeproc{ref-MouraUllrich2021}{Moura and Ullrich 2021}), using only
the core library, so that the development can be checked by anyone with
the Lean toolchain and no further dependencies. The source accompanies
the paper. Three things are verified.

First, the capacity view. Populations are represented as finite lists of
integer welfare levels, and the five conditions are stated exactly as in
Section~\ref{sec-conditions}, with ``\(A \subset W_x\), \(N(A) = n\)''
rendered as a list of \(n\) copies of \(x\) and ``\(B \subset R(1,y)\)''
as the hypothesis that every member of \(B\) lies between 1 and \(y\).
The theorem \texttt{sixth\_theorem\_countermodel} states that for every
\(A \geq 4\) and every baseline \(c \geq 0\) the view is complete and
transitive and satisfies \texttt{EgalitarianDominance},
\texttt{GeneralNonExtremePriority}, \texttt{NonElitism},
\texttt{WeakNonSadism} and \texttt{WeakQualityAddition}, and that it
does not satisfy \texttt{WeakQualityAdditionStar}. The proof is the case
analysis of Section~\ref{sec-capacity}, mechanized; it is about 250
lines.

Second, Thomas's reconstruction and the repaired theorem. Here
populations are represented as count functions from levels to numbers of
lives, with finite support, which makes the reordering of unions in the
chain proofs trivial; conditions are stated in Thomas's forms, with
``very high'' fixed at \(A\) and ``very low positive'' at \(\{1,2,3\}\),
for an arbitrary reflexive transitive relation on distributions. The
theorems \texttt{T1} through \texttt{T5} and \texttt{T6star} verify that
Thomas's six theorems hold as he states them, the last being that
Egalitarian Dominance, General Non-Extreme Priority, Non-Elitism and
Weak Non-Sadism entail VRA*; \texttt{prop2} and \texttt{cor3} verify
Proposition 2 and Corollary 3 in the same vocabulary. The two lemmas on
which the last three rest, that General Non-Extreme Priority entails
Sufficient Tradeoffs and that (General) Non-Elitism entails (General)
Inequality-Averse Addition, are proved by the double inductions that
Thomas's recursive definitions of the witnesses require. Thomas's
statements and proofs are correct in every case; the formalization
needed no repairs beyond making explicit the finiteness and range
side-conditions that his prose leaves implicit.

Thomas's conditions differ from Arrhenius's exact ones in both
directions: his General Non-Extreme Priority fixes the high and low
thresholds where Arrhenius lets them be chosen for each level, and takes
the low lives to be at a single level where Arrhenius allows a mixture.
So the theorem \texttt{T6star\_exact} repeats the argument for the exact
conditions, with the thresholds chosen inside the conditions and the low
population given as an arbitrary list of levels in the range, and
derives the negation of Weak Quality Addition* in its exact form. Its
proof differs from Thomas's in two places: the lives raised to very high
welfare in different applications of General Non-Extreme Priority must
be sent to a common level, chosen above every threshold the condition
returns, and the excellent lives in the background and in the comparison
population may sit at two different levels, which are merged by an extra
application of Inequality-Averse Addition before the main one.

Third, fidelity. The formal statements follow Arrhenius's text, with
three deliberate choices recorded in the source. Arrhenius requires a
welfare range to have at least three levels, so the range \(R(1,y)\)
carries the side condition \(y \geq 3\), which his prose statements
leave implicit. In Egalitarian Dominance the common population size is
required to be positive, since with empty populations the condition
would say that the empty population is better than itself. And in Weak
Non-Sadism, where Arrhenius writes only ``a number of lives,'' the
countermodel is shown to satisfy the condition with a positive number of
very negative lives, which is the stronger claim, while the theorems
assume only the weaker form in which that number may be zero. Since the
countermodel is verified against the exact conditions and the corrected
theorem is verified against both Thomas's and the exact conditions,
there is no tension between a verified counterexample and a verified
theorem: the counterexample is to the original conclusion, and the
theorem has the corrected one.

All proofs are complete, contain no \texttt{sorry}, and depend only on
the three standard axioms of Lean's kernel (propositional
extensionality, quotient soundness, and choice). The whole development
is about twelve hundred lines. The graded variant of
Section~\ref{sec-repair} and the solver of Section~\ref{sec-method} are
not part of it.

\section{Conclusion}\label{sec-conclusion}

Arrhenius's sixth impossibility theorem is false. The capacity view,
which ranks populations first by whether their excellent lives cover
their misery and then by total welfare, satisfies all five of its
conditions in their exact formulations, on integer-indexed welfare
levels, while accepting the Repugnant Conclusion and rejecting the Very
Repugnant Conclusion. The theorem becomes true when the Weak Quality
Addition Condition is strengthened so that the uncompensable quantity of
misery is fixed before the background population, as Thomas and Thornley
have proposed, and the strengthened theorem has been machine-checked
along with Thomas's reconstructions of the other five. But the
strengthening is not free. For every axiology on which misery can be
compensated by lives barely worth living, Weak Quality Addition in
either form is equivalent to the Quality Addition Condition that the
sixth theorem was meant to weaken; and for the axiologies on which it
cannot, the strengthened condition rules out, by stipulation, the view
that how much misery a population can absorb depends on how much
excellence it contains. What the four remaining conditions actually
establish is that the Repugnant Conclusion holds relative to any
background that has redeemed its misery. Whether that is repugnant is a
question for people who have made up their minds about the Repugnant
Conclusion itself; it is not a new result about misery.

The broader lesson concerns how results of this kind should be checked.
The error in the sixth theorem survived publication in two venues and
five years of citation before Thomas found it, and it has survived a
further decade since, because finding a gap in a proof and finding out
whether the claim is true are different tasks, and nobody undertook the
second. It was a quantifier dependence that a proof assistant rejects in
seconds, and the countermodel is one that a systematic search for
lexical potentials locates in minutes. Population axiology has
accumulated a stock of impossibility theorems whose premises are each
the product of careful philosophical judgment and whose proofs are each
a page of inequalities that almost nobody re-derives. The judgment
cannot be automated. The inequalities can, and should be.

\protect\phantomsection\label{refs}
\begin{CSLReferences}{1}{0}
\bibitem[\citeproctext]{ref-Arrhenius2000b}
Arrhenius, Gustaf. 2000a. {``An Impossibility Theorem for Welfarist
Axiologies.''} \emph{Economics and Philosophy} 16 (2): 247--66. doi:
\href{https://doi.org/10.1017/S0266267100000249}{10.1017/S0266267100000249}.

\bibitem[\citeproctext]{ref-Arrhenius2000a}
---------. 2000b. \emph{Future Generations: A Challenge for Moral
Theory}. Uppsala: Uppsala University.

\bibitem[\citeproctext]{ref-Arrhenius2003}
---------. 2003. {``The Very Repugnant Conclusion.''} In \emph{Logic,
Law, Morality: Thirteen Essays in Practical Philosophy in Honour of
Lennart {Å}qvist}, edited by Krister Segerberg and Rysiek Sliwinski.
Vol. 51. Uppsala Philosophical Studies. Uppsala: Uppsala University.

\bibitem[\citeproctext]{ref-Arrhenius2009}
---------. 2009. {``One More Axiological Impossibility Theorem.''} In
\emph{Logic, Ethics, and All That Jazz: Essays in Honour of Jordan
Howard Sobel}, edited by Lars-Göran Johansson, Jan Österberg, and Rysiek
Sliwinski. Vol. 57. Uppsala Philosophical Studies. Uppsala: Uppsala
University.

\bibitem[\citeproctext]{ref-Arrhenius2011}
---------. 2011. {``The Impossibility of a Satisfactory Population
Ethics.''} In \emph{Descriptive and Normative Approaches to Human
Behavior}, edited by Ehtibar N. Dzhafarov and Lacey Perry, 3:1--26.
Advanced Series on Mathematical Psychology. Singapore: World Scientific.
doi:
\href{https://doi.org/10.1142/9789814368018_0001}{10.1142/9789814368018\_0001}.

\bibitem[\citeproctext]{ref-ArrheniusForthcoming}
---------. Forthcoming. \emph{Population Ethics: The Challenge of Future
Generations}. Oxford: Oxford University Press.

\bibitem[\citeproctext]{ref-ArrheniusBudolfsonSpears2021}
Arrhenius, Gustaf, Mark Budolfson, and Dean Spears. 2021. {``Does
Climate Change Policy Depend Importantly on Population Ethics?''} In
\emph{Philosophy and Climate Change}, edited by Mark Budolfson, Tristram
McPherson, and David Plunkett, 111--36. Oxford: Oxford University Press.
doi:
\href{https://doi.org/10.1093/oso/9780198796282.003.0006}{10.1093/oso/9780198796282.003.0006}.

\bibitem[\citeproctext]{ref-ArrheniusEtAl2022}
Arrhenius, Gustaf, Krister Bykvist, Tim Campbell, and Elizabeth
Finneron-Burns, eds. 2022. \emph{The Oxford Handbook of Population
Ethics}. New York: Oxford University Press. doi:
\href{https://doi.org/10.1093/oxfordhb/9780190907686.001.0001}{10.1093/oxfordhb/9780190907686.001.0001}.

\bibitem[\citeproctext]{ref-BlackorbyBossertDonaldson2005}
Blackorby, Charles, Walter Bossert, and David Donaldson. 2005.
\emph{Population Issues in Social Choice Theory, Welfare Economics, and
Ethics}. Cambridge: Cambridge University Press. doi:
\href{https://doi.org/10.1017/CCOL0521825512}{10.1017/CCOL0521825512}.

\bibitem[\citeproctext]{ref-BlackorbyDonaldson1991}
Blackorby, Charles, and David Donaldson. 1991. {``Normative Population
Theory: A Comment.''} \emph{Social Choice and Welfare} 8 (3). doi:
\href{https://doi.org/10.1007/BF00177664}{10.1007/BF00177664}.

\bibitem[\citeproctext]{ref-BrandtGeist2016}
Brandt, Felix, and Christian Geist. 2016. {``Finding Strategyproof
Social Choice Functions via {SAT} Solving.''} \emph{Journal of
Artificial Intelligence Research} 55: 565--602. doi:
\href{https://doi.org/10.1613/jair.4959}{10.1613/jair.4959}.

\bibitem[\citeproctext]{ref-Broome2004}
Broome, John. 2004. \emph{Weighing Lives}. Oxford: Oxford University
Press. doi:
\href{https://doi.org/10.1093/019924376X.001.0001}{10.1093/019924376X.001.0001}.

\bibitem[\citeproctext]{ref-Carlson2022}
Carlson, Erik. 2022. {``On Some Impossibility Theorems in Population
Ethics.''} In \emph{The Oxford Handbook of Population Ethics}, edited by
Gustaf Arrhenius, Krister Bykvist, Tim Campbell, and Elizabeth
Finneron-Burns, 204--25. New York: Oxford University Press. doi:
\href{https://doi.org/10.1093/oxfordhb/9780190907686.013.14}{10.1093/oxfordhb/9780190907686.013.14}.

\bibitem[\citeproctext]{ref-GeistEndriss2011}
Geist, Christian, and Ulle Endriss. 2011. {``Automated Search for
Impossibility Theorems in Social Choice Theory: Ranking Sets of
Objects.''} \emph{Journal of Artificial Intelligence Research} 40:
143--74. doi:
\href{https://doi.org/10.1613/jair.3126}{10.1613/jair.3126}.

\bibitem[\citeproctext]{ref-Greaves2017}
Greaves, Hilary. 2017. {``Population Axiology.''} \emph{Philosophy
Compass} 12 (11): e12442. doi:
\href{https://doi.org/10.1111/phc3.12442}{10.1111/phc3.12442}.

\bibitem[\citeproctext]{ref-Huemer2008}
Huemer, Michael. 2008. {``In Defence of Repugnance.''} \emph{Mind} 117
(468): 899--933. doi:
\href{https://doi.org/10.1093/mind/fzn079}{10.1093/mind/fzn079}.

\bibitem[\citeproctext]{ref-MouraUllrich2021}
Moura, Leonardo de, and Sebastian Ullrich. 2021. {``The {L}ean 4 Theorem
Prover and Programming Language.''} In \emph{Automated Deduction -- CADE
28}, edited by André Platzer and Geoff Sutcliffe, 12699:625--35. Lecture
Notes in Computer Science. Cham: Springer. doi:
\href{https://doi.org/10.1007/978-3-030-79876-5_37}{10.1007/978-3-030-79876-5\_37}.

\bibitem[\citeproctext]{ref-Ng1989}
Ng, Yew-Kwang. 1989. {``What Should We Do about Future Generations?
{I}mpossibility of {P}arfit's {T}heory {X}.''} \emph{Economics and
Philosophy} 5 (2): 235--53. doi:
\href{https://doi.org/10.1017/S0266267100002406}{10.1017/S0266267100002406}.

\bibitem[\citeproctext]{ref-Nipkow2009}
Nipkow, Tobias. 2009. {``Social Choice Theory in {HOL}: {A}rrow and
{G}ibbard-{S}atterthwaite.''} \emph{Journal of Automated Reasoning} 43
(3): 289--304. doi:
\href{https://doi.org/10.1007/s10817-009-9147-4}{10.1007/s10817-009-9147-4}.

\bibitem[\citeproctext]{ref-ParentBenzmuller2024}
Parent, Xavier, and Christoph Benzmüller. 2024. {``Conditional Normative
Reasoning as a Fragment of {HOL}.''} \emph{Journal of Applied
Non-Classical Logics} 34 (4): 561--92. doi:
\href{https://doi.org/10.1080/11663081.2024.2386917}{10.1080/11663081.2024.2386917}.

\bibitem[\citeproctext]{ref-Parfit1984}
Parfit, Derek. 1984. \emph{Reasons and Persons}. Oxford: Oxford
University Press.

\bibitem[\citeproctext]{ref-Parfit2016}
---------. 2016. {``Can We Avoid the Repugnant Conclusion?''}
\emph{Theoria} 82 (2): 110--27. doi:
\href{https://doi.org/10.1111/theo.12097}{10.1111/theo.12097}.

\bibitem[\citeproctext]{ref-SpearsBudolfson2021}
Spears, Dean, and Mark Budolfson. 2021. {``Repugnant Conclusions.''}
\emph{Social Choice and Welfare} 57 (3): 567--88. doi:
\href{https://doi.org/10.1007/s00355-021-01321-2}{10.1007/s00355-021-01321-2}.

\bibitem[\citeproctext]{ref-TangLin2009}
Tang, Pingzhong, and Fangzhen Lin. 2009. {``Computer-Aided Proofs of
{A}rrow's and Other Impossibility Theorems.''} \emph{Artificial
Intelligence} 173 (11): 1041--53. doi:
\href{https://doi.org/10.1016/j.artint.2009.02.005}{10.1016/j.artint.2009.02.005}.

\bibitem[\citeproctext]{ref-Tannsjo2002}
Tännsjö, Torbjörn. 2002. {``Why We Ought to Accept the Repugnant
Conclusion.''} \emph{Utilitas} 14 (3): 339--59. doi:
\href{https://doi.org/10.1017/S0953820800003642}{10.1017/S0953820800003642}.

\bibitem[\citeproctext]{ref-Thomas2016}
Thomas, Teruji. 2016. {``Reconstructing {A}rrhenius's Impossibility
Theorems.''} Unpublished note.
\url{https://users.ox.ac.uk/~mert2060/webfiles/Reconstructing-Arrhenius-for-web.pdf}.

\bibitem[\citeproctext]{ref-Thomas2018}
---------. 2018. {``Some Possibilities in Population Axiology.''}
\emph{Mind} 127 (507): 807--32. doi:
\href{https://doi.org/10.1093/mind/fzx047}{10.1093/mind/fzx047}.

\bibitem[\citeproctext]{ref-Thornley2021}
Thornley, Elliott. 2021. {``The Impossibility of a Satisfactory
Population Prospect Axiology (Independently of {F}inite
{F}ine-{G}rainedness).''} \emph{Philosophical Studies} 178 (11):
3671--95. doi:
\href{https://doi.org/10.1007/s11098-021-01621-4}{10.1007/s11098-021-01621-4}.

\bibitem[\citeproctext]{ref-Thornley2022}
---------. 2022. {``A Dilemma for Lexical and {A}rchimedean Views in
Population Axiology.''} \emph{Economics and Philosophy} 38 (3):
395--415. doi:
\href{https://doi.org/10.1017/S0266267121000213}{10.1017/S0266267121000213}.

\bibitem[\citeproctext]{ref-ZuberEtAl2021}
Zuber, Stéphane, Nikhil Venkatesh, Torbjörn Tännsjö, Christian Tarsney,
H. Orri Stefánsson, Katie Steele, Dean Spears, et al. 2021. {``What
Should We Agree on about the Repugnant Conclusion?''} \emph{Utilitas} 33
(4): 379--83. doi:
\href{https://doi.org/10.1017/S095382082100011X}{10.1017/S095382082100011X}.

\end{CSLReferences}




\end{document}
